\section{Plan de trabajo}

Para el desarrollo de este trabajo de investigación se propone el siguiente plan de trabajo, dividido en etapas con sus respectivas actividades y cronograma, el cual se desarrollará durante un período de dos semestres académicos y se detallará en su respectiva sección más adelante:

\subsection{Revisión documental}

Durante esta etapa, se realizará la búsqueda y recopilación de datasets de audio de guitarra acústica, con la mayor cantidad de información adicional posible como notaciones o archivos MIDI para mejorar la calidad de las etiquetas y garantizar la mejor calidad de los datos posible.

\subsection{Implementación del modelo}
En esta fase, se llevará a cabo la implementación de los modelos de aprendizaje automático seleccionados para la detección de acordes, empleando prácticas de codificación eficientes y asegurando la reproducibilidad de los resultados. Se documentará el proceso de implementación y se prepararán los scripts necesarios para el entrenamiento y evaluación de los modelos.

\subsection{Preprocesamiento de datos}

En esta fase, se llevará a cabo la limpieza y normalización de los datos del o los datasets elegidos para el proyecto, así como la extracción de características relevantes para la detección de acordes. Se evaluarán diferentes técnicas de preprocesamiento para mejorar la calidad de las representaciones de audio.

\subsection{Entrenamiento del modelo}

En esta etapa, se implementarán y entrenarán diferentes modelos de aprendizaje automático para la detección de acordes, utilizando las características extraídas en la fase anterior. Se optimizarán los hiperparámetros de los modelos para mejorar su rendimiento.

\subsection{Comparación de modelos}

Una vez entrenados los modelos, s.e realizará una comparación exhaustiva de su desempeño utilizando métricas de evaluación como exact‑match accuracy, top‑3 accuracy y macro‑F1. Se analizarán los resultados para determinar cuál modelo ofrece el mejor rendimiento en la tarea de detección de acordes.

\subsection{Pruebas}

Finalmente, se evaluará la robustez del mejor modelo frente a variaciones de afinación, técnicas de ejecución y presencia de ruido de fondo, comparando los resultados obtenidos con una revisión manual bajo los criterios de teoría musical. Se documentarán los resultados obtenidos y se discutirán las limitaciones del sistema propuesto, así como posibles mejoras y aplicaciones futuras en el ámbito de la transcripción musical automática y el análisis de audio.

\newpage

\section{Cronograma}

El cronograma propuesto para el desarrollo de la investigación está distribuido en un período de ocho (8) meses a partir de la aprobación del anteproyecto. El Cuadro \ref{tab:cronograma} ilustra la distribución temporal de las fases del plan de trabajo.

\begin{table}[H]
\centering
\caption{Cronograma de actividades del proyecto}
\label{tab:cronograma}
\begin{tabular}{|l|c|c|c|c|c|c|c|c|}
\hline
\textbf{Actividad} & \textbf{M1} & \textbf{M2} & \textbf{M3} & \textbf{M4} & \textbf{M5} & \textbf{M6} & \textbf{M7} & \textbf{M8} \\
\hline
Revisión documental & $\bullet$ & $\bullet$ & & & & & & \\
\hline
Preprocesamiento de datos & & $\bullet$ & $\bullet$ & & & & & \\
\hline
Implementación del modelo & & & & $\bullet$ & $\bullet$ & & & \\
\hline
Entrenamiento del modelo & & & & & $\bullet$ & $\bullet$ & $\bullet$ & \\
\hline
Comparación de modelos & & & & & & & $\bullet$ & \\
\hline
Pruebas & & & & & & & & $\bullet$ \\
\hline
\end{tabular}
\end{table}

Al finalizar cada una de estas actividades, se espera cumplir con los siguientes entregables principales:

\begin{itemize}
\item \textbf{Mes 1 - 2 | Revisión documental:} Datasets seleccionados y documentados, junto con la revisión bibliográfica consolidada.
\item \textbf{Mes 2 - 3 | Preprocesamiento de datos:} Pipeline de código funcional con la extracción de características espectrales validadas.
\item \textbf{Mes 4 - 5 | Implementación del modelo:} Arquitecturas base implementadas y listas para su entrenamiento.
\item \textbf{Mes 5 - 7 | Entrenamiento del modelo:} Modelos entrenados con hiperparámetros optimizados y resultados preliminares registrados.
\item \textbf{Mes 7 | Comparación de modelos:} Reporte de desempeño comparativo con las métricas establecidas (exact-match accuracy, top-3 accuracy, macro-F1).
\item \textbf{Mes 8 | Pruebas:} Validación de robustez finalizada y entrega del documento de tesis.
\end{itemize}